
\documentclass[12pt]{article}
\usepackage[english]{babel}
\usepackage[utf8x]{inputenc}
\usepackage{amsmath}
\usepackage{graphicx}
\usepackage[font=small,labelfont=bf]{caption}
\usepackage[colorinlistoftodos]{todonotes}
\usepackage{csquotes}
\usepackage[official]{eurosym}
\usepackage{fancyhdr}
\usepackage[a4paper,bindingoffset=0.2in,%
            left=1in,right=1in,top=1in,bottom=1in,%
            footskip=.25in]{geometry}
\usepackage{blindtext}
\usepackage{minted}

\newgeometry{left=0.8in,right=0.8in,top=1in,bottom=1in}
\pagestyle {fancy}
\fancyhf{}
\rhead{Mobile Computing Project Report}
\lhead{\leftmark}
\rfoot{\thepage}

\begin{document}
\begin{titlepage}

\newcommand{\HRule}{\rule{\linewidth}{0.25mm}} % Defines a new command for the horizontal lines, change thickness here

\center % Center everything on the page

%----------------------------------------------------------------------------------------
%	HEADING SECTIONS
%----------------------------------------------------------------------------------------

\textsc{\LARGE Faculdade de Engenharia da Universidade do Porto}\\[0.5cm] % Name of your university/college
\textsc{\Large Mestrado Integrado em Engenharia Informática e Computação}\\[0.5cm] % Major heading such as course name
\textsc{\large Mobile Computing}\\[0.5cm] % Minor heading such as course title

%----------------------------------------------------------------------------------------
%	TITLE SECTION
%----------------------------------------------------------------------------------------

\HRule \\[0.75cm]
{ \huge \bfseries Mobile Computing Project Report}\\[0.4cm] % Title of your document
\HRule \\[1cm]

%----------------------------------------------------------------------------------------
%	AUTHOR SECTION
%----------------------------------------------------------------------------------------

\begin{minipage}{0.6\textwidth}
\begin{flushleft} \large
\emph{Autores:}\\
Diogo \textsc{Reis} - up201505472\\
Tiago \textsc{Magalhães} - up201607931\\

\end{flushleft}
\end{minipage}
~
\begin{minipage}{0.35\textwidth}
\begin{flushright} \large
\end{flushright}
\end{minipage}\\[1cm]

%----------------------------------------------------------------------------------------
%	LOGO SECTION
%----------------------------------------------------------------------------------------
\includegraphics[width=50mm,scale=0.5]{feuplogo.png}\\[0.5cm] % FEUP Logo



% If you don't want a supervisor, uncomment the two lines below and remove the section above
%\Large \emph{Author:}\\
%John \textsc{Smith}\\[3cm] % Your name

%----------------------------------------------------------------------------------------
%	DATE SECTION
%----------------------------------------------------------------------------------------

{\large November 2019}\\[2cm] % Date, change the \today to a set date if you want to be precise

%----------------------------------------------------------------------------------------

\vfill % Fill the rest of the page with whitespace

\end{titlepage}
\tableofcontents
\pagebreak

\section{Architecture}

\subsection{Server}

\subsubsection{REST WebAPI}
\hspace{0.6cm}
The server uses a REST WebAPI provided by the ASP.NET framework to handle all the communication with it. All routes are present in a single controller under the router of server. 

\subsubsection{Database}
\hspace{0.6cm}
Communication with the PostgreSQL database is done using the Npgsql package, it is then abstracted into a singleton class that provides the common insert and select operation. Prepared
statements are done via an entry system where each entry has a name, a value and a boolean indicating if the entry is a UUID or not, the entry names are then matched to the
query markers for parameters.

All responses to select queries are done via list of dictionaries that map strings to object, each element in the list is a row in the database, and the elements of the dictionaries
contain the values of the various columns in the database, the entries have the same name as they do in the database schema.

\subsubsection{RSA Encryption \& Signing}
\hspace{0.6cm}
Encryption is done using the .NET Crypto Service Providers, specifically RSA Crypto Service Provider, as well as the BouncyCastle package that is used for parsing and exporting public and private
keys to and from the PEM string format. This is also done with a singleton class that exposes Encryption, Decryption, Signing and Verifying data methods.

Any encrypted data or signatures are in Base64 format, while all data to be verified or encrypted are expected to be in a Unicode UTF-16 encoding.

\subsection{Store And Client Applications}

\subsubsection{HTTP Requests}
\hspace{0.6cm}
All HTTP requests are done via static methods that make use of standard Java facilities such as URL and HttpUrlConnection. All requests are run in an Android AsyncTask so as not
to block the UI main thread of Android. Each request method also receive an instant of a class called HTTPResultHandler, this class handles the result of the HTTP request once
it is ready, and this is done via its Handler method.

\subsubsection{QR Codes}
\hspace{0.6cm}
QR Code handling is done via the Google Zebra Crossing (zxing) library and some of its android integrations, methods are provided for generating a
QR code and for reading a QR code.

QR code generation is rather straightforward with a method call that returns a Bitmap.

QR code reading is a bit more involved as it makes use of the external bar code scanner application and requires the activity to override its onActivityResult method, therefore
to simplify this system and force QR code reader users to properly overload this method we have created an Android Activity that is abstract and forces its subclasses to
implement a handler that it will internally call when the onActivityResult method is triggered.

\subsubsection{RSA Encryption \& Signing}
\hspace{0.6cm}
RSA encryption and signing is done via the standard java facilities for security purposes and Androids KeyStore facilities. When a new user is registed a new private key public key
pair is introduced into the Android Secure Key Store Enclave.

All further usages of this class are done via static methods that when needed receive the username of the user, so that they can fetch his key for various security purposes.

This class offers static methods for encryption, decryption, signing and verifying data.

\subsubsection{Data Caching}
\hspace{0.6cm}
All data caching is done via the Android SharedPreferences Facilities.


\pagebreak
\section{Data Scheme}

\subsection{Database Data Schema}
\begin{minted}{postgresql}
CREATE EXTENSION IF NOT EXISTS "uuid-ossp";

DROP TABLE IF EXISTS Client CASCADE;
CREATE TABLE Client(
	id uuid primary key default uuid_generate_v4(),
	name text not null,
	username text not null,
	password text not null,
	credit_card int not null,
	public_key text not null,
	current_total_spent_euro INTEGER not null default 0,
	current_total_spent_cent INTEGER not null default 0,
	current_accumulated_euro INTEGER not null default 0,
	current_accumulated_cent INTEGER not null default 0
);

DROP TABLE IF EXISTS Voucher CASCADE;
CREATE TABLE Voucher(
	id uuid primary key default uuid_generate_v4(),
	client uuid not null REFERENCES Client(id),
	was_used BOOLEAN not null DEFAULT FALSE
);

DROP TABLE IF EXISTS Purchase CASCADE;
CREATE TABLE Purchase(
	id uuid primary key default uuid_generate_v4(),
	client uuid not null REFERENCES Client(id),
	voucher uuid REFERENCES Voucher(id) DEFAULT NULL,
	should_discount BOOLEAN not null DEFAULT false
);

DROP TABLE IF EXISTS Product CASCADE;
CREATE TABLE Product(
	id uuid primary key default uuid_generate_v4(),
	price_euro INTEGER not null,
	price_cent INTEGER not null,
	name text not null,
	image_url text DEFAULT NULL,
	purchase uuid REFERENCES Purchase(id) DEFAULT null
);
\end{minted}
\subsubsection{Client}
\hspace{0.6cm}
This table contains information about users of the platform and has the following fields:
\begin{itemize}
    \item id - UUID representing the identification of the user.
    \item name - String representing the name of the user.
    \item username - String representing the username or nickname of the user.
    \item password - String representing the password of the user.
    \item credit\textunderscore card - Integer representing the credit card number of the user.
    \item public\textunderscore key - String containing the users RSA public key.
    \item current\textunderscore total\textunderscore spent\textunderscore euro - The amount of money spent by the user, this is the euro component.
    \item current\textunderscore total\textunderscore spent\textunderscore cent - The amount of money spent by the user, this is the cent component.
    \item current\textunderscore accumulated\textunderscore euro - The amount of money the user has accumulated from voucher, this is the euro component.
    \item current\textunderscore accumulated\textunderscore cent - The amount of money the user has accumulated from voucher, this is the cent component.
\end{itemize}


\subsubsection{Voucher}
\hspace{0.6cm}
This table contains the voucher registered in the system and has the following details:
\begin{itemize}
    \item id - UUID representing the identification of the voucher.
	\item client - UUID representing the identification of the user the voucher belongs to.
	\item was\textunderscore used - Boolean representing wether or not the voucher has been used, by default this value is false.
\end{itemize}


\subsubsection{Purchase}
\hspace{0.6cm}
This table contains information about purchases and has the following details:
\begin{itemize}
    \item id - UUID representing the identification of the purchase.
	\item client - UUID representing the identification of the user the purchase is associated with.
	\item voucher - UUID representing the identification of the voucher used in this purchase, this value is optional and is null by default.
	\item should\textunderscore discount - Boolean representing weather or not the cost of the purchase was amortized with money the user had accumulated via vouchers.
\end{itemize}


\subsubsection{Product}
\hspace{0.6cm}
This table contains the products registered in the system and has the following information:
\begin{itemize}
    \item id - UUID representing the identification of the product.
	\item price\textunderscore euro - This is the price of the product, this is the euros component.
	\item price\textunderscore cent - This is the price of the product, this is the cents component.
	\item name - String representing the name of the product.
	\item image\textunderscore url - String representing the link to an image of the product.
	\item purchase - UUID representing the purchase this product is in, this is optional and is null by default.
\end{itemize}


\subsection{Checkout Information}

\subsection{Data Verification Signatures}

\pagebreak
\section{Features \& Tests}

\subsection{Features}

\subsubsection{Available Product List In Store App}

\subsubsection{User Registering}

\subsubsection{User Login}

\subsubsection{User Login Caching}

\subsubsection{Adding Products To Cart Through QR Codes}

\subsubsection{Viewing and Editing Cart}

\subsubsection{Viewing Past Transaction}

\subsubsection{Viewing Products Bough In Past Transaction}

\subsubsection{Viewing Vouchers Available To User And Amount Of Money Usable For Discounting}

\subsubsection{Checkout}

\subsection{Testing}

\pagebreak
\section{Usage Manual}

\subsection{Server}

\subsection{Store Application}
\subsubsection{Viewing Available Product List}

\subsubsection{Checkout}

\subsection{Client Application}
\subsubsection{User Registering}

\subsubsection{User Login}

\subsubsection{User Login Caching}

\subsubsection{Adding Products To Cart Through QR Codes}

\subsubsection{Viewing and Editing Cart}

\subsubsection{Viewing Past Transaction}

\subsubsection{Viewing Products Bough In Past Transaction}

\subsubsection{Viewing Vouchers Available To User And Amount Of Money Usable For Discounting}

\subsubsection{Checkout}

\end{document}
